\documentclass[11pt]{article}
\usepackage{amsmath, amssymb, amsthm}
\usepackage[margin=1in]{geometry}
\usepackage{setspace}
\usepackage{booktabs}
\setstretch{1.15}

\newtheorem{result}{Result}[section]
\newtheorem{assumption}{Assumption}

\title{Trade Tariffs and Exchange Rates in a Multilateral World\\ \large Theoretical derivations}
\author{}
\date{\today}

\begin{document}
\maketitle

\section{Introduction}

In a two-country model, an import tariff appreciates the tariffing country's currency under the Marshall--Lerner condition. In a multilateral setting, a discriminatory tariff also reallocates expenditure across foreign suppliers, so bilateral exchange-rate responses need not share a sign. These notes characterize the distinction in an $N$-country model. The general equilibrium response is $(-J)^{-1}F$, where $F$ is the vector of trade-balance disturbances created by the tariff at unchanged exchange rates and $J$ is the exchange-rate adjustment mechanism. Under symmetry, the tariffing country always appreciates against the target; its bilateral rate against an untreated bystander reverses iff cross-origin substitution exceeds a threshold $N\rho_1^*$; and its nominal effective exchange rate appreciates for every elasticity and every $N$, because the cross-origin terms cancel from the aggregate numerator. The exact cancellation and the linear threshold are properties of the symmetric nested-CES benchmark; the disturbance-adjustment decomposition, and the distinction between aggregate home-versus-foreign adjustment and reallocation across foreign suppliers, are more general. Closed forms are verified numerically in the replication package.

\section{$N$-Country Model and Equilibrium}

\subsection{Production and preferences}

Countries $i, j \in \{1, \dots, N\}$; $A$ is the tariff imposer, $B$ the target, $C$ a bystander. Each country produces a country-specific tradable and a non-tradable from labor:
\begin{equation}
Y_{T_i} = A_{T_i} L_{T_i}, \qquad Y_{N_i} = A_{N_i} L_{N_i}, \qquad L_{T_i} + L_{N_i} = L_i ,
\end{equation}
with perfect competition and producer-price normalization $P_{T_i} = 1$ in $i$'s currency, so
\begin{equation}
w_i = A_{T_i}, \qquad P_{N_i} = A_{T_i}/A_{N_i} .
\label{eq:wage}
\end{equation}
Wages are constant in own currency; all relative adjustment runs through exchange rates. Preferences are nested CES with share-linear weights:
\begin{align}
C_i &= C_{T_i}^{\alpha_T}\, C_{N_i}^{1-\alpha_T}
&& \text{(tradable share } \alpha_T)
\label{eq:outer}\\[2pt]
C_{T_i} &= \Big[\, \alpha_D^{1/\eta}\, C_{T_i i}^{\frac{\eta-1}{\eta}}
+ (1-\alpha_D)^{1/\eta}\, M_i^{\frac{\eta-1}{\eta}} \Big]^{\frac{\eta}{\eta-1}}
&& \text{(home weight } \alpha_D, \text{ elasticity } \eta)
\label{eq:mid}\\[2pt]
M_i &= \Big[\, \textstyle\sum_{j \neq i} \beta_{ij}^{1/\rho}\, C_{T_j i}^{\frac{\rho-1}{\rho}} \Big]^{\frac{\rho}{\rho-1}}
&& \text{(origin weights } \beta_{ij}, \text{ elasticity } \rho) .
\label{eq:inner}
\end{align}
$\eta$ governs substitution between domestic and imported tradables (the home-versus-foreign margin); $\rho$ governs substitution across foreign origins conditional on importing (the foreign-supplier reallocation margin), inactive at $N = 2$ and active at $N \geq 3$. The two margins are economically more general than nested CES; the constant-elasticity representation and the exact closed forms below are model-specific. The \emph{symmetric benchmark} imposes $\beta_{ij} = 1/(N-1)$, common parameters $(\alpha_T, \alpha_D, \eta, \rho)$, and equal endowments and productivities; home bias $\alpha_D$ keeps home and foreign tradables asymmetric ($\rho = \eta$, $\alpha_D = 1/N$ is not imposed).

\subsection{Prices, expenditure, income, and trade balance}

Tariffs $\tau_{ij} \geq 0$ are ad valorem, levied by $i$ on imports from $j$, with $\tau_{ii} = 0$; consumer prices in $i$'s currency are $p_{ij} = (1+\tau_{ij})\, e_{ij}$, $p_{ii} = 1$. Price indices:
\begin{equation}
P_{M_i} = \Big[ \textstyle\sum_{j\neq i} \beta_{ij}\, p_{ij}^{\,1-\rho} \Big]^{\frac{1}{1-\rho}},
\qquad
P_{T_i} = \Big[ \alpha_D\, p_{ii}^{\,1-\eta} + (1-\alpha_D)\, P_{M_i}^{\,1-\eta} \Big]^{\frac{1}{1-\eta}},
\qquad
P_i = \kappa_i\, P_{T_i}^{\,\alpha_T}\, P_{N_i}^{\,1-\alpha_T} .
\label{eq:indices}
\end{equation}
Expenditure shares of income:
\begin{equation}
s_{ii} = \alpha_T\, \alpha_D \left(\frac{p_{ii}}{P_{T_i}}\right)^{1-\eta},
\qquad
s_{ij} = \alpha_T\, (1-\alpha_D) \left(\frac{P_{M_i}}{P_{T_i}}\right)^{1-\eta}
\beta_{ij} \left(\frac{p_{ij}}{P_{M_i}}\right)^{1-\rho}, \quad j \neq i ,
\label{eq:sfor}
\end{equation}
tariff-inclusive: $s_{ij} I_i$ is spending at consumer prices, $s_{ij} I_i/(1+\tau_{ij})$ the border value. Revenue is rebated lump-sum,
\begin{equation}
I_i = w_i L_i + \sum_{j \neq i} \frac{\tau_{ij}}{1+\tau_{ij}}\, s_{ij} I_i
\quad\Longrightarrow\quad
I_i = \frac{w_i L_i}{1 - \sum_{j\neq i} \frac{\tau_{ij}}{1+\tau_{ij}}\, s_{ij}} .
\label{eq:income}
\end{equation}
Trade balances at producer prices in a common currency (the tariff wedge is a domestic transfer):
\begin{equation}
\mathrm{TB}_i = \sum_{k \neq i} f_{ki} - \sum_{j \neq i} f_{ij},
\qquad
f_{ij} \equiv \frac{s_{ij}\, I_i}{1+\tau_{ij}} .
\label{eq:tb}
\end{equation}
By Walras' law only $N-1$ conditions are independent. Closure: no internationally traded assets, $\mathrm{TB}_i = 0$ for all $i$; given $\{\tau_{ij}\}$, the $N-1$ free rates solve $\mathrm{TB}_i = 0$, $i \neq A$.

\subsection{Exchange-rate concepts}

$e_{ij}$ = units of $i$'s currency per unit of $j$'s, $e_{ii} = 1$, $e_{ij} = e_{ik}e_{kj}$; \textbf{a rise in $e_{ij}$ is a depreciation of $i$ against $j$}. Perturbations $\nu_j \equiv d\log e_{Aj}$, $\nu_A = 0$. Definitions:
\begin{equation}
\omega_{ij} \equiv \frac{e_{ij} w_j}{w_i},
\qquad
\mathrm{ToT}_{ij} \equiv \frac{P_{T_i}}{e_{ij} P_{T_j}} = \frac{1}{e_{ij}},
\qquad
q_{ij} \equiv \frac{e_{ij} P_j}{P_i},
\label{eq:rer}
\end{equation}
\begin{equation}
d\nu_i^{E} \equiv \sum_{j \neq i} w_{ij}\, d\log e_{ij},
\qquad
d q_i^{E} \equiv \sum_{j \neq i} w_{ij}\, d\log q_{ij},
\qquad
w_{ij} = \frac{f_{ij} + f_{ji}}{\sum_{k\neq i} (f_{ik} + f_{ki})} ,
\label{eq:neer}
\end{equation}
with $w_{ij} = 1/(N-1)$ in the symmetric model; effective rates are depreciation-positive. Under the production normalization,
\begin{equation}
d\log e_{ij} = d\log \omega_{ij} = -\,d\log \mathrm{ToT}_{ij} ,
\label{eq:onetoone}
\end{equation}
so every nominal exchange-rate result is equivalently a relative-wage or terms-of-trade result.

\subsection{General linearization}

Linearize at a balanced free-trade baseline with normalized incomes: home share $h_i = \alpha_D$ (a country-specific $h_i$ carries through unchanged), within-import shares $b_{ij} = \beta_{ij}$, income shares $s_{ij} = m\, b_{ij}$ with $m \equiv \alpha_T(1-\alpha_D)$, flows $f_{ij} = s_{ij} I_i$, $\sum_k f_{ki} = \sum_j f_{ij}$. Perturbing by tariffs $dt_{ij}$:
\begin{align}
dp_{ij} &= \nu_j - \nu_i + dt_{ij} \qquad (dp_{ii} = 0),
\label{eq:dp}\\
d\log P_{M_i} &= \sum_{j\neq i} b_{ij}\, dp_{ij},
\qquad
d\log P_{T_i} = (1-\alpha_D)\, d\log P_{M_i} ,
\label{eq:dindex}\\
d\log s_{ii} &= -(1-\eta)(1-\alpha_D)\, d\log P_{M_i},
\label{eq:dshome}\\
d\log s_{ij} &= \underbrace{(1-\rho)\big( dp_{ij} - d\log P_{M_i} \big)}_{\text{across origins}}
+ \underbrace{(1-\eta)\,\alpha_D\, d\log P_{M_i}}_{\text{home vs.\ imports}} ,
\label{eq:dsfor}\\
d\log I_i &= \nu_i + \sum_{j\neq i} s_{ij}\, dt_{ij} ,
\label{eq:dincome}\\
d\,\mathrm{TB}_i
&= \sum_{k \neq i} f_{ki} \big( d\log s_{ki} + d\log I_k - dt_{ki} \big)
- \sum_{j \neq i} f_{ij} \big( d\log s_{ij} + d\log I_i - dt_{ij} \big) .
\label{eq:dtb}
\end{align}
Stacking $d\nu \equiv (\nu_j)_{j \neq A}$ and $d\boldsymbol{\tau} \equiv (dt_{jk})_{j \neq k}$, the conditions for $i \neq A$ are
\begin{equation}
J\, d\nu + G\, d\boldsymbol{\tau} = 0,
\qquad
J_{il} \equiv \frac{\partial\, d\mathrm{TB}_i}{\partial \nu_l},
\qquad
G_{i,(jk)} \equiv \frac{\partial\, d\mathrm{TB}_i}{\partial\, dt_{jk}}\bigg|_{\nu \text{ fixed}} .
\label{eq:JFdef}
\end{equation}
A tariff experiment is a direction $a$ scaled by $d\tau$ (unilateral tariff: $a_{AB} = 1$; war: $a_{AB} = a_{BA} = 1$), with impact vector
\begin{equation}
F \equiv G a
\qquad\Longrightarrow\qquad
J\, d\nu = -F\, d\tau,
\qquad
\frac{d\nu}{d\tau} = (-J)^{-1} F .
\label{eq:system}
\end{equation}
$F$ is the vector of trade-balance disturbances created by the tariff at unchanged exchange rates; $J$ describes how trade balances respond to exchange rates; configurations differ only in the direction $a$ applied to the same $G$. This local disturbance-adjustment representation holds for any differentiable multilateral trade environment with $N-1$ independent trade-balance conditions; nested CES supplies the particular entries above.

\subsection{The multilateral Marshall--Lerner condition}

Let $\widetilde J$ be the full $N \times N$ Jacobian before dropping country $A$. Homogeneity (only relative prices enter \eqref{eq:dp}) and Walras' law give the identities
\begin{equation}
\sum_{l} \widetilde J_{il} = 0 \quad \forall i,
\qquad\qquad
\sum_{i} \widetilde J_{il} = 0 \quad \forall l .
\label{eq:hom}
\end{equation}

\begin{assumption}[Gross substitutability]
$\widetilde J_{il} \geq 0$ for all $l \neq i$. In the symmetric systems below every off-diagonal entry is a positive multiple of $D_N$, so the assumption there is $D_N > 0$, which holds for $\eta \geq 1$.
\label{as:gs}
\end{assumption}

Homogeneity and Assumption \ref{as:gs} give $\widetilde J_{ii} = -\sum_{l\neq i}\widetilde J_{il} \leq 0$ (the bilateral Marshall--Lerner conditions), and the reduced $J$ inherits a dominant diagonal,
\begin{equation}
|J_{ii}| = \sum_{l \neq i,\,A} J_{il} + \widetilde J_{iA} \;\geq\; \sum_{l \neq i,\,A} J_{il} ,
\label{eq:diagdom}
\end{equation}
strict in every row with $\widetilde J_{iA} > 0$. By the standard characterization of irreducibly diagonally dominant Z-matrices:

\begin{result}[Multilateral Marshall--Lerner condition]
Under Assumption \ref{as:gs}, with $J$ irreducible and $\widetilde J_{iA} > 0$ for some $i$, $-J$ is a nonsingular M-matrix:
\begin{equation}
\det(-J) > 0,
\qquad
(-J)^{-1} \geq 0 \ \text{elementwise},
\qquad
\operatorname{Re}\lambda(J) < 0 .
\label{eq:invpos}
\end{equation}
\label{res:mml}
\end{result}

\begin{result}[Sign rule]
Under \eqref{eq:invpos},
\begin{equation}
\operatorname{sign}\!\left( \frac{d\nu_i}{d\tau} \right)
= \operatorname{sign}\!\Big( \sum_j \omega_{ij} F_j \Big),
\qquad \omega_{ij} \equiv \big[(-J)^{-1}\big]_{ij} \geq 0 .
\label{eq:signrule}
\end{equation}
\label{res:signrule}
\end{result}

\noindent Every exchange-rate response is a nonnegative combination of the impact disturbances, so opposite-signed bilateral responses require a mixed-sign $F$: bilateral ambiguity comes from the tariff disturbance, not from a perverse adjustment mechanism. Gross substitutability is sufficient, not necessary; the eigenvalue face of \eqref{eq:invpos} is local stability of the tatonnement $\dot\nu_i = k_i \mathrm{TB}_i$, so stability is not a separate requirement. Since $F$ is linear in $\rho$ and $\operatorname{adj}(-J)$ has degree $N-2$, each weighted sum in \eqref{eq:signrule} is a degree-$(N-1)$ polynomial in $\rho$: one relevant root at $N = 3$, possibly several sign changes at $N \geq 4$.

\section{Tariff Responses: Two, Three, and $N$ Countries}

Throughout this section $A$ imposes $d\tau = dt_{AB}$ at the symmetric baseline, and
\begin{equation}
\rho_1^*
\equiv 1 + \alpha_D(\eta-1) - \alpha_T(1-\alpha_D)
= \underbrace{\alpha_D\,\eta}_{\text{home switching}} + \underbrace{(1-\alpha_D)(1-\alpha_T)}_{\text{expenditure absorption}}
\;>\; 0 .
\label{eq:rho1-identity}
\end{equation}

\subsection{Two countries}

At $N = 2$ the system \eqref{eq:system} is scalar and $\rho$ is inoperative ($dp_{AB} = d\log P_{M_A}$). With $s_{AB} = s_{BA} = m$, equations \eqref{eq:dp}--\eqref{eq:dtb} collect to
\begin{equation}
d\,\mathrm{TB}_A
= \underbrace{m\, D_2\, \nu_B}_{\text{adjustment}}
+ \underbrace{m\, \rho_1^*\, d\tau}_{\text{impact disturbance}}
= 0,
\qquad
D_2 \equiv 2\alpha_D(\eta-1) + 1 ,
\label{eq:2c-collected}
\end{equation}
\begin{equation}
\boxed{\;
\frac{d\log e_{AB}}{d\tau}\bigg|_{\tau=0}
= -\,\frac{\rho_1^*}{D_2} .
\;}
\label{eq:2c-slope}
\end{equation}
In reduced form, with export volume $X(1/e)$, import volume $M((1+\tau)e)$, and elasticities $\varepsilon_X, \varepsilon_M$, balanced trade gives $\partial \mathrm{TB}_A/\partial e = M(\varepsilon_X + \varepsilon_M - 1)$; comparing with \eqref{eq:2c-collected} per unit of import value,
\begin{equation}
\varepsilon_X + \varepsilon_M - 1 = D_2,
\qquad
\varepsilon_X = \varepsilon_M = \alpha_D(\eta-1) + 1 .
\label{eq:ml-D2}
\end{equation}

\begin{result}[Conventional wisdom, $N=2$]
Condition \eqref{eq:invpos} collapses to $D_2 > 0$, which is exactly the textbook Marshall--Lerner condition $\varepsilon_X + \varepsilon_M > 1$ (automatic for $\eta \geq 1$). Under it a unilateral import tariff appreciates the tariffing country's currency, since $\rho_1^* > 0$.
\label{res:2c}
\end{result}

\noindent At unchanged exchange rates the tariff improves $A$'s balance by $m\rho_1^* d\tau$, through home switching and non-tradable absorption; the appreciation eliminates the surplus. The disturbance signs the numerator and Marshall--Lerner signs the denominator. (Appendix A gives the exact finite-tariff Cobb--Douglas solution, which confirms the sign globally.)

\subsection{Three countries}

At $N = 3$ ($\beta_{ij} = \tfrac12$, $s_{ij} = m/2$), \eqref{eq:dp}--\eqref{eq:dindex} give
\begin{equation}
d\log P_{M_A} = \tfrac{1}{2}(\nu_B + \nu_C + d\tau), \qquad
d\log P_{M_B} = \tfrac{1}{2}\nu_C - \nu_B, \qquad
d\log P_{M_C} = \tfrac{1}{2}\nu_B - \nu_C ,
\label{eq:3c-indices}
\end{equation}
with $d\log I_A = \tfrac{m}{2} d\tau$, $d\log I_B = \nu_B$, $d\log I_C = \nu_C$. Substituting into \eqref{eq:dtb} for $i = B, C$ and collecting:
\begin{equation}
J = -\frac{mD_3}{4}\begin{pmatrix} 2 & -1 \\ -1 & 2 \end{pmatrix},
\qquad
F = -\frac{m}{4}\begin{pmatrix} \rho + \rho_1^* \\ \rho_1^* - \rho \end{pmatrix},
\qquad
D_3 = 3\alpha_D(\eta-1) + \rho + 1 .
\label{eq:3c-JF}
\end{equation}
The impact disturbances are $F_B = -\tfrac{m}{4}(\rho + \rho_1^*) < 0$ and $F_C = \tfrac{m}{4}(\rho - \rho_1^*)$: the target always loses; the bystander gains diverted expenditure (increasing in $\rho$) against reduced sales to a poorer $B$ and $A$'s home switching, with net gain iff $\rho > \rho_1^*$. Inverting:
\begin{equation}
\begin{pmatrix} \nu_B \\ \nu_C \end{pmatrix}
= (-J)^{-1} F\, d\tau
= -\frac{1}{3D_3}
\begin{pmatrix} 2 & 1 \\ 1 & 2 \end{pmatrix}
\begin{pmatrix} \rho + \rho_1^* \\ \rho_1^* - \rho \end{pmatrix} d\tau
= \frac{1}{3D_3}
\begin{pmatrix} -(\rho + 3\rho_1^*) \\ \rho - 3\rho_1^* \end{pmatrix} d\tau .
\label{eq:3c-solution}
\end{equation}

\begin{result}[Three-country threshold]
With $D_3 > 0$,
\begin{equation}
\frac{d\log e_{AB}}{d\tau} = -\,\frac{\rho + 3\rho_1^*}{3D_3} < 0,
\qquad
\boxed{\;\frac{d\log e_{AC}}{d\tau} = \frac{\rho - 3\rho_1^*}{3D_3}\;}
\qquad
\frac{d\log e_{BC}}{d\tau} = \frac{2\rho}{3D_3} > 0 :
\label{eq:3c-slopes}
\end{equation}
$A$ always appreciates against $B$, $B$ always depreciates against $C$, and $A$ depreciates against the untreated $C$ iff $\rho > \rho_3^* \equiv 3\rho_1^*$.
\label{res:3c}
\end{result}

\noindent The gap between the impact threshold $\rho_1^*$ and the equilibrium threshold $3\rho_1^*$ is the sign rule at work: $d\nu_C/d\tau \propto F_B + 2F_C \propto \rho - 3\rho_1^*$, so $C$'s equilibrium rate mixes its own diversion gain with its exposure to $B$'s deterioration. A bystander's favorable impact disturbance need not determine its equilibrium bilateral rate, because that rate responds to the entire vector $F$ through $(-J)^{-1}$.

\emph{Corollary (single elasticity).} With $\rho = \eta = \sigma$: $\rho_3^* - \sigma = \sigma(3\alpha_D - 1) + 3(1-\alpha_D)(1-\alpha_T) > 0$ whenever $\alpha_D \geq 1/3$, for any $\alpha_T$; the reversal is unreachable. At the equal-weights point $\alpha_D = 1/3$, $\alpha_T = 3/4$: $d\log e_{AC}/d\tau = -1/(12\sigma) \to 0^-$.

\subsection{General symmetric $N$}

With weights $1/(N-1)$ and the $N-2$ bystanders identical (common $\nu_C$), \eqref{eq:dp}--\eqref{eq:dindex} give
\begin{equation}
d\log P_{M_A} = \frac{\nu_B + d\tau + (N-2)\nu_C}{N-1},
\qquad
d\log P_{M_B} = \frac{(N-2)\nu_C}{N-1} - \nu_B,
\qquad
d\log P_{M_C} = \frac{\nu_B - 2\nu_C}{N-1},
\label{eq:N-indices}
\end{equation}
with $d\log I_A = \tfrac{m}{N-1}d\tau$, $d\log I_B = \nu_B$, $d\log I_C = \nu_C$. Collecting \eqref{eq:dtb} for $B$ and a representative bystander (rows scaled by $(N-1)/m$; verified symbolically):
\begin{equation}
D_N
\begin{pmatrix} -(N-1) & N-2 \\ 1 & -2 \end{pmatrix}
\begin{pmatrix} \nu_B \\ \nu_C \end{pmatrix}
=
\begin{pmatrix} (N-2)\rho + \rho_1^* \\ \rho_1^* - \rho \end{pmatrix} d\tau ,
\qquad
D_N = N\alpha_D(\eta-1) + (N-2)\rho + 1 ,
\label{eq:N-system}
\end{equation}
where the counting matrix has determinant $N$; inverting:
\begin{equation}
\frac{d\log e_{AB}}{d\tau}
= -\,\frac{N\rho_1^* + (N-2)\rho}{N\, D_N} < 0,
\qquad
\frac{d\log e_{AC}}{d\tau}
= \frac{\rho - N\rho_1^*}{N\, D_N},
\qquad
\frac{d\log e_{BC}}{d\tau}
= \frac{(N-1)\,\rho}{N\, D_N} > 0 ,
\label{eq:N-slopes}
\end{equation}
which reduce to \eqref{eq:3c-slopes} at $N = 3$ and to \eqref{eq:2c-slope} at $N = 2$.

\begin{result}[Dilution threshold]
The bystander bilateral rate reverses iff
\begin{equation}
\boxed{\;\rho > \rho_N^* = N \rho_1^* . \;}
\end{equation}
\label{res:Nthresh}
\end{result}

\noindent Diverted expenditure is spread across $N-2$ bystanders, each with weight $1/(N-1)$, so stronger cross-origin substitution is required to reverse any individual bilateral rate (dilution). The exact linearity in $N$ is a symmetric nested-CES result. With $\rho = \eta = \sigma$, $\rho_N^* - \sigma = \sigma(N\alpha_D - 1) + N(1-\alpha_D)(1-\alpha_T) > 0$ whenever $\alpha_D \geq 1/N$: the reversal stays unreachable for single-elasticity models on the home-biased region.

\section{Bilateral versus Effective Exchange Rates}

Averaging \eqref{eq:N-slopes} with symmetric weights:
\begin{equation}
\frac{d\nu_A^E}{d\tau}
= \frac{\nu_B + (N-2)\nu_C}{N-1}
= \boxed{\; -\,\frac{\rho_1^*}{D_N} \;} < 0,
\qquad
\nu_B - \nu_C = -\,\frac{(N-1)\rho}{N D_N}\, d\tau .
\label{eq:N-neer}
\end{equation}

\begin{result}[NEER preservation]
Under symmetry, a unilateral tariff appreciates the tariffing country's NEER for every $N \geq 2$ and every $\rho$ satisfying \eqref{eq:invpos}, even when $\rho > N\rho_1^*$ reverses the bilateral rate against every bystander.
\label{res:neer}
\end{result}

\noindent $\rho$ governs the relative distribution of adjustment across foreign partners: its direct terms cancel from the NEER numerator under symmetry, while it remains in $D_N$ and therefore still affects the magnitude of the aggregate response. The two-country appreciation result thus survives multilateralization as an effective, rather than universally bilateral, statement --- it is bilateral at $N = 2$ only because there is one foreign country. The exact cancellation is specific to the symmetric nested-CES environment (Section 5.3 quantifies its robustness); the distinction between aggregate home-versus-foreign adjustment and foreign-supplier reallocation is more general. (Appendix B collects the correspondence across $N$ and remarks on the empirical mapping of $\rho$.)

\section{Extensions}

\subsection{Alternative tariff configurations}

Responses are linear in the tariff vector, so configurations are superpositions of directed bilateral shocks (directions $a$ on the same $G$).

\paragraph{Broad unilateral tariff.} A uniform tariff on all $N-1$ partners treats every origin identically, deactivating cross-origin reallocation. Superposing \eqref{eq:N-slopes} over targets:
\begin{equation}
\frac{d\log e_{Aj}}{d\tau} = -\,\frac{(N-1)\rho_1^*}{D_N} = \frac{d\nu_A^E}{d\tau}
\qquad \forall j :
\label{eq:N-broad}
\end{equation}
$\rho$ drops out of the numerator, every bilateral rate appreciates equally, and bilateral and effective responses coincide. Bilateral ambiguity arises from discrimination across foreign suppliers, not protection itself.

\paragraph{Symmetric trade war.} Adding the mirror shock $dt_{BA} = d\tau$:
\begin{equation}
\frac{d\log e_{AB}}{d\tau}\bigg|^{w} = 0,
\qquad
\frac{d\log e_{AC}}{d\tau}\bigg|^{w} = \frac{\rho - \rho_1^*}{D_N},
\qquad
\frac{d\nu_A^{E}}{d\tau}\bigg|^{w} = \frac{N-2}{N-1}\,\frac{\rho - \rho_1^*}{D_N} .
\label{eq:N-war}
\end{equation}

\begin{result}[War threshold]
In a symmetric war the belligerents' effective rates depreciate, and every bystander appreciates against both, iff
$\rho > \rho_1^*$, independently of $N$.
\label{res:war}
\end{result}

\noindent Retaliation cancels the direct bilateral component between $A$ and $B$, leaving the bystander reallocation margin decisive: each bystander's equation decouples, so its own impact effect sets the sign and does not dilute with $N$. The war threshold is far below the unilateral threshold $N\rho_1^*$ because the direct appreciation against $B$ has been removed. At $\eta = 1$, $\rho_1^* = 1 - m < 1$.

\subsection{Real exchange rates}

From \eqref{eq:rer} and \eqref{eq:indices}, $d\log q_{Aj} = \nu_j + \alpha_T(d\log P_{T_j}^{\mathrm{own}} - d\log P_{T_A}^{\mathrm{own}})$; for the unilateral tariff both bilateral real rates take one form at every $N$ (verified symbolically):
\begin{equation}
d\log q_{Aj} = \Big( 1 - \frac{mN}{N-1} \Big)\, \nu_j \;-\; \frac{m}{N-1}\, d\tau,
\qquad j = B, C :
\label{eq:qgen}
\end{equation}
attenuated nominal pass-through plus a direct price-index term working toward real appreciation of $A$. Setting $d\log q_{AC} = 0$ with \eqref{eq:N-slopes} (for $m < 1/N$):
\begin{equation}
\rho_q^*(N) = \frac{\big[(N-1) - mN\big]\, N\rho_1^* + mN\big[ N\alpha_D(\eta-1) + 1 \big]}{(N-1)(1 - mN)}
\;>\; N\rho_1^* ,
\label{eq:N-rhoq}
\end{equation}
with benchmark values $\rho_q^*(3, 4, 5) = 4.93,\, 7.34,\, 10.40$ against nominal thresholds $3.96,\, 5.28,\, 6.60$. Averaging \eqref{eq:qgen}:
\begin{equation}
\frac{dq_A^E}{d\tau}
= -\Big( 1 - \frac{mN}{N-1} \Big) \frac{\rho_1^*}{D_N} - \frac{m}{N-1}
\;<\; 0
\qquad \text{for all } \rho, N \ \big(m < \tfrac{N-1}{N}\big) .
\label{eq:N-reer}
\end{equation}

\begin{result}[Real rates]
The direct CPI effect makes real bilateral reversal harder than nominal reversal, $\rho_q^*(N) > N\rho_1^*$, and reinforces effective appreciation: the REER response is negative for all $\rho$ and $N$, and strictly larger in magnitude than the pass-through of the nominal response alone.
\label{res:real}
\end{result}

\subsection{Asymmetry}

The $J$--$F$ framework remains valid without symmetry; the exact thresholds and the NEER cancellation do not.

\paragraph{General asymmetric baseline.} At any balanced baseline, \eqref{eq:dp}--\eqref{eq:dtb} hold with observed shares $(b_{ij}, h_i)$ and incomes $y_i$ as coefficients. For $N = 3$ at a general baseline (paper appendix), the bystander slope has denominator $\det(-J) > 0$ and numerator quadratic in $\rho$,
\begin{equation}
\operatorname{sign} \frac{d\log e_{AC}}{d\tau}
= \operatorname{sign}\big( c_2\, \rho^2 + c_1\, \rho + c_0 \big),
\qquad
c_2 = b_{AB}\, b_{AC}\, b_{CA}\, b_{CB}\, (1-h_A)(1-h_C)\, y_A\, y_C ,
\label{eq:asym-quad}
\end{equation}
with $c_1, c_0$ depending on the full pattern of shares and sizes. Since $c_2 > 0$ whenever $A$ and $C$ each buy from both origins, sufficiently strong cross-origin substitution reverses $e_{AC}$ at any such baseline; the threshold is the largest root, equal to $3\rho_1^*$ at the symmetric point. Away from symmetry no universal bilateral threshold exists.

\paragraph{Structured asymmetry.} Let the target carry a multiple $\kappa$ of the symmetric bilateral share (comparable across $N$, since the symmetric share shrinks like $1/(N-1)$), with the weight matrix symmetric and balanced:
\begin{equation}
b_{AB} = b_{BA} = \frac{\kappa}{N-1},
\qquad
b_{AC} = b_{CA} = b_{BC} = b_{CB} = \frac{1 - \frac{\kappa}{N-1}}{N-2},
\qquad
b_{CC'} = \frac{1 - 2\,b_{CA}}{N-3} .
\label{eq:kappa-b}
\end{equation}
The reduced system again has two unknowns and solves in closed form (symbolically; replication package). With NEER weights $w_{Aj} = b_{Aj}$, the response is a rational function of $(\rho, N, \kappa)$ collapsing to $-\rho_1^*/D_N$ at $\kappa = 1$, with numerator $-\kappa\rho_1^*$ times a quadratic $Q(\rho)$ whose leading coefficient is proportional to $(\kappa-1)(N-1-\kappa)$.

\begin{result}[Asymmetric NEER]
Holding $\kappa$ fixed as $N \to \infty$,
\begin{equation}
N\,\frac{d\nu_A^E}{d\tau}
\;\longrightarrow\;
-\,\kappa\;\frac{\rho_1^*}{\rho + \alpha_D(\eta-1)}
= -\,\kappa \lim_{N\to\infty}\frac{N\rho_1^*}{D_N} .
\label{eq:kappa-neer}
\end{equation}
\label{res:kappa}
\end{result}

\noindent To leading order, asymmetric exposure scales the symmetric NEER response by the target's relative importance and preserves its sign. At finite $N$: for $\kappa \geq 1$ all coefficients of $Q$ are positive at large $N$, so effective appreciation holds for every $\rho$; for $\kappa < 1$, $Q$ has a single positive root $\rho_E^*(\kappa, N) = \tfrac{N\rho_1^*}{1-\kappa}(1 + O(N^{-1}))$, beyond which the NEER depreciates --- an under-weighted target lets depreciations against many bystanders outweigh the appreciation against $B$, with $\rho_E^* \to N\rho_1^*$ as $\kappa \to 0$. The reversal threshold is of order $N$ ($\rho_E^* \approx 22$ at $N = 10$, $\kappa = 0.5$, benchmark). The symmetric NEER result is exact under symmetry and approximately robust within this structured asymmetric family; robustness to arbitrary asymmetric trade networks is not claimed.

\section{Relation to the Canonical PTA Model}

\subsection{Competing exporters}

Importer $A$ holds only the numeraire $y$; exporters $B, C$ hold good $x$; preferences are quasilinear, $U_i = u_i(x_i) + y_i$. $A$ levies specific tariffs $t_B, t_C$; with $q_j$ the exporter price and $p$ the domestic price,
\begin{equation}
p = q_B + t_B = q_C + t_C,
\qquad
M_A(p) = E_B(p - t_B) + E_C(p - t_C),
\qquad M_A' < 0, \; E_j' > 0 ,
\label{eq:pta-mc}
\end{equation}
where $M_A(p) = x_A^d(p)$ and $E_j(q_j) = \omega_j - x_j^d(q_j)$ from $u_i' = $ local price. Differentiating in $t_B$:
\begin{equation}
\frac{dp}{dt_B} = \frac{E_B'}{E_B' + E_C' - M_A'} \in (0, 1),
\qquad
\frac{dq_B}{dt_B} = -\,\frac{E_C' - M_A'}{E_B' + E_C' - M_A'} < 0,
\qquad
\frac{dq_C}{dt_B} = \frac{dp}{dt_B} > 0 .
\label{eq:pta-cs}
\end{equation}

\begin{result}[Viner--Mundell]
A discriminatory tariff on $B$ worsens $B$'s and improves $C$'s terms of trade, unconditionally.
\label{res:viner}
\end{result}

\subsection{Relation to the present model}

An improvement in $C$'s terms of trade is a rise in $e_{AC}$; the nested model requires $\rho > 3\rho_1^*$, the canonical model nothing. The canonical environment removes the margins that oppose diversion toward $C$. Decomposing the threshold (benchmark $\alpha_D = 0.8$, $\alpha_T = 0.4$, $\eta = 1.5$):
\begin{equation}
\rho_3^* = 3\rho_1^* = \underbrace{3}_{\text{$B$--$C$ linkage}}
+ \underbrace{3\alpha_D(\eta-1)}_{\text{home switching}}
- \underbrace{3\alpha_T(1-\alpha_D)}_{\text{openness credit}}
= 3.0 + 1.2 - 0.24 = 3.96 .
\label{eq:pta-decomp}
\end{equation}
The canonical model deletes the first two opposing forces and the income effects behind the first: $B$ and $C$ are competing exporters rather than full trading partners (no $B$--$C$ demand linkage), $A$ does not produce the imported good (no home switching), and quasilinearity removes income effects. What remains is reallocation between the two foreign suppliers, signed directly by $M_A' < 0$ and $E_j' > 0$; the Viner case adds $\rho = \infty$.

\begin{result}
The canonical PTA model is the restriction of the three-country model to a subspace on which $\rho_3^* \leq 0$: the unconditional Viner--Mundell sign is the threshold condition evaluated where it cannot bind. The nested model embeds the same diversion mechanism among the opposing margins, and $\rho > \rho_3^*$ states when it dominates.\footnote{Analytical PTA models obtain sharp signs by such restrictions (competing exporters: Bagwell--Staiger 1999, Ornelas 2005; symmetric endowment blocs: Krugman 1991, Bond--Syropoulos 1996; Viner: $\rho = \infty$; Kemp--Wan 1976: no restrictions, no sign claims). Applied GTAP-class and quantitative trade models retain every margin and generate both signs in simulations without stating the separating condition. Numerically, severing the $B$--$C$ linkage alone accounts for most of the threshold, and combining it with $\alpha_D = 0$ removes the reversal region entirely (replication package).}
\label{res:pta}
\end{result}

\section{Conclusion}

The general equilibrium response to a tariff is $(-J)^{-1}F$: the tariff creates trade-balance disturbances at unchanged exchange rates, and the adjustment mechanism, regular under the multilateral Marshall--Lerner condition, combines them with nonnegative weights. At $N = 2$ this reduces to the textbook Marshall--Lerner appreciation result. At $N \geq 3$, a discriminatory tariff generates foreign-supplier reallocation and hence a mixed-sign disturbance vector: the tariffing country always appreciates against the target, but under symmetry depreciates against a bystander iff $\rho > N\rho_1^*$, a threshold that rises with $N$ because diversion is diluted across bystanders. Aggregation removes the direct reallocation term from the NEER numerator, so the effective exchange rate appreciates for every $\rho$ and $N$; configuration and asymmetry set the limits of this aggregate result, with reciprocal retaliation reversing it at $\rho > \rho_1^*$ and an under-weighted target reversing it at $\rho$ of order $N$.

The organizing distinction is between aggregate home-versus-foreign adjustment, which determines the sign of the effective response, and reallocation across foreign suppliers, which distributes the adjustment across bilateral rates and affects the magnitude of the aggregate response. Tariffs change both how much a country buys from abroad and from whom it buys it.

\appendix

\section{Exact Cobb--Douglas solution ($N = 2$)}

At $\eta = 1$ shares are constant at the weights, $I_A = w_A L_A/[1 - \tfrac{\tau}{1+\tau}m]$ from \eqref{eq:income}, and $\mathrm{TB}_A = 0$ reads $e\, m\, w_B L_B = m I_A/(1+\tau)$:
\begin{equation}
e(\tau) = \frac{w_A L_A}{w_B L_B}\cdot \frac{1}{1 + (1-m)\tau},
\qquad
\frac{d\log e}{d\tau}\bigg|_{\tau=0} = -(1-m) = -\frac{\rho_1^*}{D_2}\bigg|_{\eta=1} ,
\label{eq:2c-exact}
\end{equation}
globally decreasing in $\tau$: the linearized sign in \eqref{eq:2c-slope} is exact at all tariff levels in this case. ($\eta = 1$ fixes the elasticity, not the weight $\alpha_D$.)

\section{Correspondence across $N$ and empirical remarks}

\begin{center}
\begin{tabular}{p{3.6cm}p{3.0cm}p{3.4cm}p{4.2cm}}
\toprule
object & $N = 2$ & $N = 3$ & general $N$ \\
\midrule
ML condition \eqref{eq:invpos} & $\varepsilon_X + \varepsilon_M - 1 = D_2 > 0$ & $D_3 > 0$ & $-J$ a nonsingular M-matrix \\[3pt]
sufficient primitive & $\eta \geq 1$ & $\eta \geq 1$ & gross substitutability (Assumption \ref{as:gs}) \\[3pt]
sign rule & one impact, unconditional & weights $\tfrac{1}{3}\left(\begin{smallmatrix} 2 & 1 \\ 1 & 2\end{smallmatrix}\right)$ & $\operatorname{sign}\big(\sum_j \omega_{ij}F_j\big)$, $\omega \geq 0$ \\[3pt]
bilateral threshold & none (no bystander) & $3\rho_1^*$ & $N\rho_1^*$ (symmetric); degree-$(N-1)$ roots in general \\[3pt]
real bilateral threshold & none & $\rho_q^*(3)$ & $\rho_q^*(N)$ \eqref{eq:N-rhoq} \\[3pt]
NEER (unilateral) & $-\rho_1^*/D_2$ & $-\rho_1^*/D_3$ & $-\rho_1^*/D_N$: sign invariant under symmetry \\[3pt]
war threshold & --- & $\rho_1^*$ & $\rho_1^*$, independent of $N$ \\
\bottomrule
\end{tabular}
\end{center}

\paragraph{Empirical mapping of $\rho$.} The parameter $\rho$ is substitution across foreign origins conditional on the import composite; estimates labeled ``trade elasticities'' generally measure different objects (aggregate import responses, home-versus-foreign substitution, gravity responses mixing margins). Qualitatively, upper-nest estimates are low and lower-nest estimates several times larger, so the benchmark $\rho_3^* = 3.96$ is neither clearly above nor below plausible $\rho$: neither sign of $e_{AC}$ can be ruled out without estimates that map onto $\rho$ directly. The war threshold $\rho_1^* = 1.32$ is less sensitive to this mapping, so the war predictions are comparatively robust.

\end{document}
